\title{Byte Latent Transformer: Patches Scale Better Than Tokens}

\begin{abstract}
We present the Byte Latent Transformer ({\textsc BLT}), a novel large language model architecture operating at the byte level that, for the first time, achieves parity with tokenization-based models in large-scale settings, while substantially enhancing inference efficiency and robustness.
{\textsc BLT} processes input bytes by grouping them into dynamically determined patches, which become the fundamental units for computation. These patches are delineated according to the entropy of subsequent bytes, thus directing greater computational and modeling resources to segments of higher data complexity. We conduct the inaugural flop-controlled scaling analysis of byte-level architectures, scaling up to 8B parameters and 4T bytes of training data. Our findings confirm that it is viable to scale models directly on raw byte sequences without requiring a predefined vocabulary. Training and inference benefit from selecting longer patches in predictable regions of data, which also leads to notable qualitative improvements in reasoning capabilities and generalization to rare cases. In summary, at a constant inference cost, {\textsc BLT} achieves more favorable scaling compared to models reliant on tokenization, by enabling concurrent increases in both patch and model size.
\end{abstract}

\section{Introduction}
\label{section:intro}

We present the Byte Latent Transformer~(\textbf{BLT}), a novel architecture that eliminates the need for tokenization by directly processing raw byte sequences. For the first time, BLT achieves parity with traditional tokenization-dependent models at scale while also offering notable gains in both computational efficiency and robustness (\S\ref{sec:robustness}). While contemporary large language models (llms) are largely trained in an end-to-end fashion, they universally rely on tokenization—a heuristic that fragments byte sequences into a predetermined set of tokens prior to model input. This tokenization process inherently shapes the compression of input strings, introducing limitations such as heightened sensitivity to particular domains or modalities~\citep{dagangetting}, vulnerability to input noise~(\S\ref{sec:robustness}), limited representation of orthographic variations~\citep{edman2024cute}, and unfair performance across different languages~\citep{liang2023xlm,petrov2024language,limisiewicz2024myte}.

Tokenization has historically played a crucial role, as training llms directly on byte sequences incurs prohibitive computational expenses at scale due to the increased sequence lengths~\citep{xue2022byt5}. Previous research addresses this issue through the use of self-attention mechanisms that are more computationally efficient~\citep{el2020characterbert, clark2022canine} or through architectures that eliminate attention altogether~\citep{wang2024mambabyte}~(\S\ref{section:related_work}). Nonetheless, such strategies predominantly benefit the training of \textit{small models}. When scaling up, the primary computational burden in Transformer models stems from the expansive feed-forward network layers that are applied to every byte, rather than from the attention components themselves.

To optimize compute usage, we introduce a flexible, adaptive approach for aggregating bytes into \textit{patches}~(\S\ref{section:patching}), complemented by a novel model architecture that integrates both byte-level and patch-level data. Distinct from tokenization, {\textsc BLT} does not rely on a predetermined set of patches. Instead, any collection of bytes can be transformed into latent patch representations through efficient, learnable encoder and decoder components. Our experiments demonstrate that this strategy yields a \textit{more} effective distribution of computational resources compared to models built on tokenization.

Conventional tokenization-based llms distribute computational resources equally across all tokens. This approach prioritizes computational efficiency but may undermine performance, since the tokenization process relies on compression strategies that do not consistently align with the actual complexity of each predictive task. The foundation of our proposed architecture rests on enabling the model to selectively increase computation for challenging segments as needed. As an illustration, most word endings present straightforward, low-entropy prediction tasks, and thus do not warrant the application of a large transformer; such extensive computation is more justified for tasks like generating the initial token of a new sentence, which typically require higher cognitive effort. In the {\textsc BLT} model~(\S\ref{section:architecture}), this philosophy is realized by utilizing three transformer modules: two compact local models operating at the byte level, and a more expansive global \textit{latent transformer}~(\autoref{fig:arch_high_level}). 

The process of forming patches—collections of bytes—and, consequently, of allocating computational effort adaptively, is guided in {\textsc BLT} by partitioning the input based on the entropy observed in next-byte prediction. This segmentation strategy ensures that byte groupings share comparable levels of informational content within their local context.

We conduct the inaugural scaling analysis of byte-level models governed by computational cost (flop-controlled), training models up to 8B parameters on 4T bytes, and demonstrate that it is feasible to scale end-to-end training directly on bytes without relying on fixed-vocabulary tokenizers. Our findings reveal that {\textsc BLT} achieves parity with Llama 3 in terms of flop-controlled training performance,\footnote{The computational cost for each model is measured by tallying the total number of Floating Point OPerations (flops) performed.} while offering up to a 50\% reduction in flops required at inference~(\S\ref{section:scaling}). Additionally, we find that modeling directly at the byte level yields notable gains in capturing the distributional nuances of rare or long-tail data instances. {\textsc BLT} models exhibit greater resilience to noisy inputs compared to models utilizing tokenizers, and display improved abilities in character-level processing, as evidenced by enhanced performance in tasks assessing orthographic recognition, phonological analysis, and translation for low-resource languages~(\S\ref{sec:robustness}). Furthermore, the design of {\textsc BLT} models enables simultaneous expansion of both model size and patch size, all while adhering to a constant inference flop budget. Increasing patch size typically reduces computational demands, allowing excess compute to be diverted towards scaling the global latent transformer, since it operates less frequently. Through scaling experiments regulated by inference flop constraints~(\autoref{fig:fixed_inference_scaling}), we document markedly superior scaling dynamics relative to architectures reliant on tokenization.

To conclude, our work offers several key advances: 1) We present {\textsc BLT}, a byte-level latent LLM framework capable of adaptively allocating computational resources to enhance flop efficiency; 2) We provide evidence that our models reach flop-controlled training equivalence with Llama 3 up to the 8B parameter range, with the added flexibility of sacrificing small amounts of performance on evaluation benchmarks to obtain flop savings of as much as 50\%; 3) The {\textsc BLT} approach introduces a novel axis for LLM scaling, permitting increases in model size while adhering to a constant inference budget; 4) We illustrate that {\textsc BLT} exhibits heightened resilience to noisy input and demonstrates a sensitivity to sub-word structure in data that token-based LLMs typically overlook. The training and inference implementations for {\textsc BLT} are publicly available at~\url{https://github.com/facebookresearch/blt}.

\section{Patching: From Individual Bytes to Groups of Bytes}
\label{section:patching}

Dividing bytes into \textit{patches} enables {\textsc BLT} to flexibly distribute computational resources according to contextual requirements. Figure~\ref{fig:patching_types} illustrates multiple approaches for breaking a byte sequence into patches. More precisely, a patching function $f_p$ partitions a byte sequence $\pmb{x}=\{x_i,|i=1,\ldots n\}$ of length $n$ into $m < n$ patches $\pmb{p}=\{p_j|j=1,\ldots,m\}$ by assigning each $x_i$ to either 0 or 1, where 1 designates the initiation of a new patch. In both patch-based and token-based frameworks, the main determinant of computational expense is the number of main Transformer operations, which, in the case of {\textsc BLT}, corresponds to the total number of patches defined by a chosen patching function. Therefore, the average patch length, or simply the \textit{patch size}, becomes the key metric for evaluating computational requirements in both training and inference phases for a specified patching function~(\S\ref{section:flops}).

In the following sections, we present three specific patching strategies: segmenting into patches of fixed byte length~(\S\ref{section:static-patch}), patching at whitespace boundaries~(\S\ref{section:white-patch}), and adaptive patching using entropy values derived from a compact byte-level language model~(\S\ref{section:dyn-patch}). We conclude by exploring incremental patching and outlining distinctions between patching and standard tokenization methods~(\S\ref{section:bpe-patch}).

\subsection{Strided Patching Every K Bytes}
\label{section:static-patch}

A simple method for aggregating bytes is to partition them into patches of a fixed length $k$, as demonstrated in MegaByte~\citep{yu2023megabyte}. Employing a constant stride simplifies both the training and inference processes, offers a direct way to modify the average patch size, and facilitates easy regulation of the computation cost. Despite these advantages, this patching strategy introduces notable drawbacks. Primarily, computational resources are not allocated adaptively based on the informational content: transformer step $j$ might be inefficiently used for segments containing mostly whitespace, or might not provide enough processing for highly informative segments such as mathematical expressions. Additionally, this method results in inconsistent and non-adaptive segmentation of similar byte strings, causing instances like the same word to be partitioned differently.

\subsection{Space Patching}
\label{section:white-patch}

\citet{slagle2024spacebyte} introduces a straightforward yet impactful enhancement to strided patching, whereby new patches are initiated after encountering any space-like bytes\footnote{
    Space-like bytes refer to any byte that does not correspond to a Latin letter, numeric digit, or UTF-8 continuation byte. Moreover, it is required that each patch includes at least one byte that is not space-like. }, which often coincide with natural linguistic boundaries in a variety of languages. In the Space patching technique, each word is explicitly modeled with a latent transformer operation (i.e., increased computational effort), dedicating additional computational resources to every word. This approach guarantees consistent patching of words across different sequences and ensures computational intensity is focused on making challenging predictions that typically appear after spaces. For instance, determining the first byte in a response to a question like ``Who composed the Magic Flute?\uline{\hspace{.6em}}'' is substantially more difficult than subsequent bytes following the initial ``M''; the first byte constrains the set of plausible answers, thus making the completion ``Mozart'' significantly easier to predict after the initial letter. Nonetheless, space patching exhibits limitations: it struggles to adapt smoothly to all languages and domains, and, crucially, does not support flexible patch sizes. The following section presents a novel patching strategy that leverages the observation that the initial bytes of words typically present the highest prediction difficulty, while also facilitating a systematic means to adjust patch size.

\subsection{Entropy Patching: Using Next-Byte Entropies from a Small Byte LM}
\label{section:dyn-patch}

Instead of utilizing a heuristic method based on rules like whitespace detection, our strategy adopts a data-driven method for determining where next-byte prediction uncertainty is greatest. We propose \textit{entropy patching}, a technique that employs entropy measurements to establish patch boundaries.

To implement this, we fit a compact byte-level auto-regressive language model to the {\textsc BLT} training set and evaluate the entropy of predicting the subsequent byte under the language model’s probability distribution $p_{e}$ defined on the byte vocabulary $\mathcal{V}$:

\begin{equation}
H(x_i) = \sum_{v \in \mathcal{V}} p_{e}(x_i=v|\pmb{x}_{<i}) \log p_{e}(x_i=v|\pmb{x}_{<i})
\end{equation}

We explore two approaches for detecting patch boundaries using entropies $H(x_i)$. The initial method selects points whose entropy exceeds a predefined global threshold, as shown in~\autoref{fig:patching}. The alternative method locates points where entropy is significantly higher compared to the preceding value. This latter strategy may also be viewed as detecting points that disrupt the roughly monotonically decreasing pattern of entropy within a patch.

\begin{align*}
    \text{Global Constraint}\hspace{1cm}H(x_t)& > \theta_g\\
    \text{Approx. Monotonic Constraint}\hspace{1cm}H(x_t)& - H(x_{t-1}) > \theta_r
\end{align*}

Patch boundaries are determined in a streamlined preprocessing phase conducted as part of dataloading. This approach contrasts with~\citet{nawrot-etal-2023-efficient}, who utilize a classifier trained to recognize entropy-based patch boundaries. Within our experiments~(\S\ref{section:experiments}), we evaluate both strategies for differentiating low- and high-entropy bytes.

\subsection{The Byte-Pair Encoding (BPE) Tokenizer and Incremental Patching}
\label{section:incremental-patching}
\label{section:bpe-patch}

Numerous state-of-the-art llms, such as our reference model Llama 3, employ subword tokenizers like bpe~\citep{gage1994new, sennrich-etal-2016-neural}. Throughout this work, we define ``tokens'' as byte sequences selected from a \textit{finite} set established before the training process, in contrast to ``patches,'' which are sequences grouped dynamically and lack a predefined vocabulary. Importantly, one distinguishing characteristic of tokens versus patches is that token-based models do not possess explicit access to the original byte-level information.

An important advancement that {\textsc BLT} introduces over tokenization-based architectures lies in reshaping the balance between vocabulary magnitude and computational requirements. Conventional large language models experience a direct trade-off: expanding the vocabulary results in, on average, longer tokens and consequently reduces the total number of generation steps, but this expansion also necessitates a larger output dimension for the model’s final projection layer. This inherent limitation constrains tokenization strategies, making it difficult for them to considerably alter both token length and inference expenses. As an illustration, Llama 3 attains an increase in the mean token size from 3.7 to 4.4 bytes, but only by quadrupling the size of its embedding table relative to Llama 2.

During generation, {\textsc BLT} must determine at each step whether the current position in the byte sequence aligns with a patch boundary, since this choice dictates if additional computation through the Latent Transformer is required. Importantly, this determination must be made without reference to subsequent, as-yet-ungenerated portions of the sequence. Consequently, the method for segmenting the byte sequence into patches cannot rely on information from future bytes. Formally, we define a patching method $f_p$ to be incrementally patchable if the following holds:
$$f_p(\pmb{x}_{<i}) = f_p(\pmb{x})_{<i}$$
The commonly used bpe algorithm is not incrementally patchable, as the way a given prefix is segmented can vary based on its continuation; thus, it fails to satisfy the condition above.\footnote{Inserting a unique delimiter token at patch edges allows bpe to fulfill the incremental patching criterion, but this approach increases the overall length of the byte sequence.}

\section{{\textsc BLT} Architecture}
\label{section:architecture}
The {\textsc BLT} framework consists of a principal global autoregressive language model that works at the patch level, complemented by a pair of compact local models: one responsible for converting sequences of bytes into patch representations, and the other for reconstructing bytes from these patch representations~(\autoref{fig:arch_high_level}).

\subsection{Latent Global Transformer Model}

\textit{The Latent Global Transformer} is an autoregressive transformer architecture, denoted as $\mathcal{G}$, consisting of $l_{\mathcal{G}}$ layers. This model receives a sequence of latent input patch representations $p_j$ and produces a corresponding sequence of output patch representations $o_j$. In this work, we use the subscript $j$ to represent individual patches and $i$ to indicate bytes. The global model implements a block-causal attention mask~\citep{dubey2024llama}, limiting each patch’s attention scope to its current and preceding patches within the same document. Since this model accounts for the majority of floating point operations during both pre-training and inference, selecting when to activate it provides a mechanism for modulating computational resources, thereby adapting compute allocation across different input and output segments based on their respective complexity.

\subsection{Local Encoder}
\label{section:local}

\textit{The Local Encoder Model}, represented as $\mathcal{E}$, consists of a compact transformer architecture with $l_{\mathcal{E}} << l_{\mathcal{G}}$ layers. Its primary function is to swiftly transform a stream of input bytes $b_i$ into rich patch-level representations $p_j$. Distinguishing it from standard transformers, the architecture introduces a cross-attention layer subsequent to each transformer layer; this addition facilitates the aggregation of byte-wise representations into patch-wise forms~(\autoref{fig:crossattn}).

Initially, the byte sequence $b_i$ is projected into embeddings $x_i$ using a parameter matrix of size $\mathbb{R}^{256 \times h_{\mathcal{E}}}$. These initial embeddings may further incorporate supplemental features, such as hash-embeddings~(\S\ref{sec:hash-emb}), depending on the configuration. Next, multiple layers consisting of alternating transformer and cross-attention mechanisms~(\S\ref{sec:enc-cross-attn}) sequentially update the representations, culminating in patch representations $p_i$, which are subsequently input into the global transformer $\mathcal{G}$. The attention mechanism within the transformer employs a \textit{local block causal} masking strategy, restricting each byte’s attention to a sliding window of $w_{\mathcal{E}}$ prior bytes; this attention window is permitted to span dynamic patch boundaries but is confined within document boundaries. The next sections elaborate on the specifics of the embedding process and the construction of the cross-attention block.

\subsubsection{Encoder Hash n-gram Embeddings}
\label{sec:ngram-emb}
\label{sec:hash-emb}
An essential aspect of constructing resilient and expressive representations at each position $i$ involves integrating context from prior bytes. In {\textsc BLT}, this is addressed by representing both the standalone byte $b_i$ and its occurrence within a byte n-gram structure. Specifically, for every index $i$, we initially generate byte-grams

\begin{align}
    g_{i,n}=\{b_{i-n + 1},\ldots, b_i\}
\end{align}

for every byte position $i$ and for values of $n$ ranging from three to eight.\footnote{
Byte-grams with size $n$ or greater are excluded in cases where $i < n$. }

Next, we present hash $n$-gram embeddings, which assign each byte $n$-gram an index in a fixed-size embedding table $E_{n}^{hash}$ through a hash function, considering all $n$ in the set $\{3,4,5,6,7,8\}$~\citep{bai2010learning}. 
Each embedding produced in this way is subsequently combined with the byte’s original embedding, after which the sum is normalized and fed as input to the local encoder model. The overall augmented embedding is computed as follows:

\begin{align}
e_i &= x_i + \sum_{n = 3,...,8}E_{n}^{hash}(\text{Hash}(g_{i,n})) \\
\text{where, Hash}(g_{i,n}) &= \text{RollPolyHash}(g_{i,n}) \% |E_{n}^{hash}|
\end{align}

We adjust $e_i$ by dividing by the total number of $n$-gram sizes plus one, and apply RollPolyHash as specified in Appendix~\ref{appendix:rollpolyhash}. Section~\ref{section:ablations} presents an ablation study examining how varying both the $n$ values and the size of the embedding table for $n$-gram hash embeddings influences scaling law behaviors under fixed computational budgets. Beyond hash-based $n$-gram embeddings, our investigation extended to frequency-based $n$-gram embeddings; further information regarding this line of inquiry can be found in Appendix~\ref{appendix:ngrams}.

\subsubsection{Encoder Multi-Headed Cross-Attention}
\label{sec:enc-cross-attn}
Our approach is largely inspired by the input cross-attention mechanism utilized in the Perceiver architecture~\citep{jaegle2021perceiver}, with a key modification: here, the latent representations dynamically map to variable-length patch embeddings rather than a fixed collection of latent vectors~(\autoref{fig:crossattn}), and the attention mechanism is restricted solely to the bytes within each corresponding patch. Each patch $p_j$ is assigned a query vector, which is generated by first aggregating the representations of the bytes within $p_j$ (e.g., via pooling), followed by the application of a linear projection, $\mathcal{E}_{C} \in \mathbb{R}^{h_{\mathcal{E}} \times (h_{\mathcal{E}} \times U_{\mathcal{E}})}$, where $U_{\mathcal{E}}$ denotes the total number of encoder cross-attention heads. To specify this process, let $f_{\text{bytes}}(p_j)$ be the byte sequence that constitutes patch $p_j$; then, we determine

\begin{align}
P_{0,j} &= \mathcal{E}_{C}(f_\text{bytes}((p_j)), f~ \text{is a pooling function} \\
P_l &= P_{l-1} + W_o\left(\text{softmax}\left(\frac{QK^T}{\sqrt{d_k}}\right)V\right) \\
\text{where } Q_j &= W_q(P_{l-1,j}), K_i = W_k(h_{l-1,i}), V_i = W_v(h_{l-1,i}) \\
h_l &= \text{Encoder-Transformer-Layer}_l(h_{l-1})
\end{align}

In this context, $P \in \mathbb{R}^{n_p \times h_{\mathcal{G}}}$ denotes the representations for $n_p$ patches, which serve as input to the global model. These representations are initialized by aggregating the byte embeddings $e_i$ corresponding to each patch $p_j$. The matrices $W_q$, $W_k$, $W_v$, and $W_o$ serve as the learned linear transformations for queries, keys, values, and output, respectively. Keys and values are derived by projecting byte-level representations $h_i$ from the preceding layer (or $e_i$ in the initial layer). A specialized masking mechanism is implemented such that each query $Q_j$ only interacts with keys and values associated with bytes within the same patch $j$. Since multi-headed attention is employed across $Q$, $K$, and $V$, and since patch representations are usually higher-dimensional ($h_{\mathcal{G}}$) compared to $h_{\mathcal{E}}$, we represent $P_l$ using multiple attention heads of size $h_{\mathcal{E}}$ during cross-attention, and subsequently concatenate these head outputs to form the $h_{\mathcal{G}}$-dimensional patch representations. Furthermore, we apply a pre-LayerNorm to the queries, keys, and values, while omitting any positional encodings within this cross-attention module. Lastly, a residual connection encompasses the cross-attention block to facilitate information flow.

\subsection{Local Decoder} 

In parallel with the local encoder, the local decoder $\mathcal{D}$ is designed as an efficient transformer-based architecture, utilizing $l_{\mathcal{D}} << l_{\mathcal{G}}$ layers. Its function is to transform a set of global patch representations $o_j$ into a corresponding sequence of raw bytes $y_i$. The decoder generates the sequence of bytes in an autoregressive manner, conditioning each prediction on the bytes previously decoded. To accomplish this, it incorporates the hidden states output by the local encoder for the given byte sequence. The decoding process is structured as $l_{\mathcal{D}}$ repeated blocks consisting of a cross-attention layer followed by a transformer layer. Initially, the cross-attention module is responsible for mapping the patch-level representations to byte-level representations, after which the transformer component of the decoder processes this byte-wise sequence.

\subsubsection{Decoder Multi-headed Cross-Attention} 

Within the decoder's cross-attention mechanism, the functions of queries and key/value pairs are reversed; specifically, byte-representations serve as the queries while patch representations act as key/values. The starting point for the byte-representations in the cross-attention comes from the byte embeddings output by the encoder’s final layer, denoted as $h_{l_{\mathcal{E}}}$. For each subsequent layer $l$, the updated byte-representations $d_{l,i}$ are calculated as follows:

\begin{align}
D_0 &= h_{l_{\mathcal{E}}} \\
B_l &= D_{l-1} + W_o\left(\text{softmax}\left(\frac{QK^T}{\sqrt{d_k}}\right)V\right), \\
  &\text{where } Q_i = W_q(d_{l-1,i}), K_i = W_k(\mathcal{D}_C(o_j)), V_i = W_v(\mathcal{D}_C(o_j)) \\
  D_l &= \text{Decoder-Transformer-layer}_l(B_l)
\end{align}

Here, as before, $W_k$ and $W_v$ serve as the projection matrices for keys and values, respectively. These matrices are responsible for applying both a linear transformation and the splitting operation $\mathcal{D}_C$ to the final patch representations $o_j$ provided by the global model. The projection matrix $W_q$ produces queries based on byte representations $d_{l-1}$ from the preceding layer in the decoder transformer (or from $h_{l_{\mathcal{E}}}$ in the first layer), while $W_o$ serves as the output projection matrix. Consequently, $B$ has shape $\mathbb{R}^{h_{\mathcal{D}} \times n_b}$, with $n_b$ representing the total number of output bytes. The subsequent decoder representations $D_l$ are obtained by applying a decoder transformer layer to $B$, the result of the cross-attention operation. Following the procedure of the local encoder’s cross-attention, we utilize multi-head attention, adopt pre-layer normalization, omit positional embeddings, and introduce a residual connection that bypasses the cross-attention component.

\section{Experimental Setup}
\label{section:experiments}
We established carefully controlled experimental protocols to compare {\textsc BLT} and tokenization-based models, making sure that {\textsc BLT} does not benefit from potentially larger sequence contexts.

\subsection{Pre-training Datasets} In all experiments discussed in this work, the models are pre-trained on two main datasets: (1) The Llama 2 dataset~\citep{touvron2023llama}, which consists of 2 trillion tokens sourced from diverse publicly accessible datasets and subsequently subjected to cleaning and filtering processes to enhance data quality; and (2) {\textsc BLT}-1T, a novel collection encompassing 1 trillion tokens obtained from assorted public domains, inclusive of a subset from the Datacomp-LM pre-training release~\citep{li2024datacomp}. The Llama 2 dataset serves as the basis for scaling law analyses, focusing on the optimal number of tokens as indicated by~\cite{dubey2024llama}, to inform the architectural design of {\textsc BLT}. In contrast, {\textsc BLT}-1T is leveraged for end-to-end pre-training to facilitate head-to-head evaluations against Llama 3 on a suite of downstream benchmarks. It is important to note that neither dataset incorporates content originating from Meta’s proprietary platforms or services. Additionally, in all baseline tokenizer model experiments, the Llama 3 tokenizer—featuring a 128K token vocabulary—is employed, as it consistently outperformed the Llama 2 tokenizer in our evaluations.

\subsection{Entropy Model} For our experiments, the entropy model is implemented as a byte-level language model, utilizing the same training dataset as the full {\textsc BLT} model. By default, this model consists of a transformer architecture with 100M parameters, composed of 14 layers, each with a hidden size of 512. It employs a sliding window attention mechanism covering 512 bytes. All other hyperparameters are aligned with those of our local and global transformer variants. As outlined in \autoref{section:ablations}, we explored the effects of varying model size, architecture, and receptive field. Notably, when the model's receptive field is sufficiently constrained, the resulting trained entropy model is compact enough to be represented via an efficient lookup table.

\subsection{Entropy Threshold and Equalizing Context Length} In entropy-based patching models, we determine a patching threshold aimed at achieving a specific average \textit{patch size} across the pretraining data mixture. Unlike tokenization, in {\textsc BLT} the \textit{patch size} is fully adjustable, which markedly affects the resulting context window available to the model. To ensure models with larger patch sizes do not benefit from an excessive context window, we control for context length by holding the total byte count per batch constant in expectation. Consequently, for models utilizing larger patch sizes, we proportionally decrease the number of sequences per batch. Specifically, with Llama 2 data, an 8k byte context is employed, whereas for the {\textsc BLT}-1T dataset, we extend the context to an average of 16k bytes, all while preserving a constant average batch size of 16M bytes.

Although the overall batch size remains unchanged, dynamic patching techniques result in varying proportions of bytes per patch when assembling data batches. To optimize performance, our {\textsc BLT} training procedure organizes patches into batches in such a way that padding operations are unnecessary for the computationally intensive latent transformer. This strategy guarantees a uniform patch count across all batches. For training, we pad and, if needed, truncate input byte sequences to 12k and 24k bytes for Llama 2 and {\textsc BLT}-1T datasets, respectively, thereby preventing memory overloads from sequences containing exceptionally large patches.

\subsection{Entropy Model Context}
\label{section:entropy-context}
Our empirical observations indicate that entropy patching leads to increasingly larger patches, particularly in highly structured datasets such as multiple choice questions (refer to the patching behavior in an MMLU case shown in \autoref{fig:mmlu_patching}). Such tasks exhibit substantial repetitiveness, contributing to lower entropy within the model's context for the redundant segments. To address this, in our extensive {\textsc BLT}-Entropy experiment with a patch size of 4.5, we reinitialize the entropy context using new lines and apply an approximate monotonicity constraint. This adjustment helps mitigate "entropy drift," which emerges as context length changes. It is important to note that this modification solely impacts entropy computation, while our methodology for selecting the entropy threshold remains unchanged.

\subsection{FLOPs Estimation} 
\label{section:flops}

Our calculation of transformer flops closely mirrors the methodology in Chinchilla~\citep{hoffmann2022training}, which incorporates the flop counts for feed-forward components, qkvo projections within self-attention mechanisms, as well as the operations required for attention computation and output projection. However, an important distinction in our approach is the assumption that the input embedding layer leverages an efficient lookup table rather than performing dense matrix multiplication, thus rendering it a zero-flop operation. Consistent with prior literature, we approximate the computational cost of the backward pass as double that of the forward pass.

To determine the flops \textit{per byte} metric for {\textsc BLT} models, we aggregate the computational cost from multiple components: the local encoder transformer, the global latent transformer, the local decoder transformer, as well as the cross attention modules situated within both the encoder and decoder. Specifically, we calculate:

\begin{align}
    \text{FL}_{\text{{\textsc BLT}}} &= \text{Transf. FL}(h_{\mathcal{G}}, l_{\mathcal{G}}, m=n_{ctx}/n_p,V=0) / n_p\\
   &+\text{Transf. FL}(h_{\mathcal{E}}, l_{\mathcal{E}}, m=w_{\mathcal{E}}, V=0) \\
   &+ \text{Transf. FL}(h_{\mathcal{D}}, l_{\mathcal{D}}, m=w_{\mathcal{D}}, V=256) \\ 
    &+ \text{Cross Attn. FL}(h_{\mathcal{E}}, l_{\mathcal{E}}, m=n_p, r=n_p/k) \times k/n_p \\ 
   &+ \text{Cross Attn. FL}(h_{\mathcal{D}}, l_{\mathcal{D}}, m=k, r=k/n_p) 
\end{align}

Here, $n_{ctx}$ represents the byte-level sequence length, and $n_p$ indicates the patch size. The parameter $r$ denotes the ratio between the number of queries and the number of keys/values, while $k$ describes the proportion of the patch dimension to the byte dimension, reflecting how many local models are concatenated to construct a global model representation (as illustrated with $k=2$ in \autoref{fig:crossattn}). $V$ signifies the vocabulary size associated with the output projection, which is utilized solely in the local decoder. The context length, $m$, employed by the attention mechanism varies depending on whether the module processes a byte sequence or a patch sequence. We update the calculation of attention FLOPs for each respective module. The precise formulas for evaluating FLOPs in both the Transformer and cross-attention components are detailed in Appendix~\ref{section:flops-appendix}.

\subsection{Bits-Per-Byte Estimation} The metric of perplexity is inherently tied to the use of a specific tokenizer, as it evaluates the level of model uncertainty per token. However, to enable fair comparison between models operating at the byte and token granularity, we adopt the Bits-Per-Byte (BPB) metric, consistent with prior studies~\citep{xue2022byt5, yu2023megabyte, wang2024mambabyte}. BPB serves as a tokenizer-agnostic analog to perplexity. Concretely, we compute:

\begin{align}
    \text{BPB}(x)&= \frac{\mathcal{L}_{CE}(\pmb{x})}{\ln(2)\cdot n_{\text{bytes}}}
\end{align}

Here, the total cross-entropy loss over the input sequence $\pmb{x}$ quantifies model uncertainty, and this sum is scaled by both the total byte count in $\pmb{x}$ and a logarithmic constant for normalization.

\subsection{Transformer Architecture Hyperparameters} 

In all transformer layers within {\textsc BLT}—encompassing both local and global modules—our design closely aligns with the architecture utilized in Llama~3~\citep{dubey2024llama}. Specifically, the feed-forward networks employ the SwiGLU activation function~\citep{swiglu}. Rotary positional embeddings (RoPE)~\citep{RoPe}, parameterized by $\theta=500000$~\citep{xiong2024effective}, are incorporated solely within the self-attention components. For layer normalization, RMSNorm~\citep{RMSNorm} is adopted. Self-attention layers that rely on fixed-standard attention masks—such as \textit{block causal} or \textit{fixed-window block causal}—leverage Flash attention~\citep{dao2022flashattention}, with a window length set to 512 when using fixed-width attention. In contrast, cross-attention modules with dynamically generated patch-specific masks utilize Flex Attention\footnote{\url{https://pytorch.org/blog/flexattention}} to enable efficient fused execution, thereby greatly accelerating the training process.

\subsection{{\textsc BLT}-Specific Hyperparameters} In our investigation of {\textsc BLT} models, we analyze both scaling dynamics and performance on downstream tasks, evaluating a range of model sizes: 400M, 1B, 2B, 4B, and 8B. The specific architectural choices for these models are detailed in Appendix~Table~\ref{tab:arch}. For query initialization in the initial cross-attention layer of the local encoder, we employ max-pooling. Consistently across all {\textsc BLT} configurations, we utilize $500,000$ hashes with a single hash function, selecting n-gram lengths from 3 to 8. Every model is trained with a fixed learning rate of $4e-4$. This selection is informed by a hyperparameter search at the 400M and 1B scales, where a sweep from $1e-3$ to $1e-4$ revealed identical optimal learning rates for both token and {\textsc BLT} models. For studies involving scaling on Llama-2 datasets, training batch sizes follow recommendations from~\cite{dubey2024llama}, or an equivalent measured in bytes. The optimization strategy involves the AdamW optimizer~\citep{adamw}, configured with $\beta_1=0.9$ and $\beta_2=0.95$, and $\epsilon=10^{-8}$. A learning rate schedule starts with 2000 steps of linear warm-up and proceeds with a cosine decay to zero. Additionally, weight decay is set at 0.1, and a global gradient clipping value of 1.0 is enforced.

\section{Scaling Trends}
\label{section:scaling}

We provide a comprehensive analysis of scaling behaviors in byte-level models to guide the ongoing scaling of {\textsc BLT} models. Our study seeks to overcome the shortcomings of earlier investigations into byte-level models by: (a) analyzing scaling patterns under compute-optimal training conditions, (b) training equivalent 8B-parameter models with substantial datasets (reaching 1T tokens or 4T bytes) and assessing their performance on downstream applications, and (c) evaluating scaling dynamics while maintaining constant inference cost. A subsequent section will explore the distinct benefits afforded by byte-sequence modeling.

\subsection{Parameter Matched Compute Optimal Scaling Trends}
\label{section:parameter-matched-scaling}
Utilizing the Llama 2 dataset, we conduct training for multiple \textit{compute-optimal} bpe and {\textsc BLT} models across four distinct parameter scales, from 1B up to 8B. For each model, we chart the number of training flops against language modeling accuracy on a representative sample of the training data mixture. The bpe models are optimized using the ideal parameter-to-data ratio, as established in Llama~3~\citep{dubey2024llama}. This \textit{compute-optimal} configuration is theoretically expected to maximize training dataset performance within the constraints of a fixed training budget~\citep{hoffmann2022training}, thereby serving as a strong reference for comparison. Alongside each bpe model, an equivalent {\textsc BLT} model is trained on identical data, employing a Latent Transformer with an architecture and parameter count that aligns with the matched bpe Transformer.

As shown in~\autoref{fig:scaling-compute-opt} (right), {\textsc BLT} models consistently equal or surpass the performance of bpe models, and this pattern persists as both model size and computational requirements increase. To our knowledge, {\textsc BLT} represents the first byte-level Transformer architecture to demonstrate comparable scaling behavior to BPE-based models within compute-optimal settings. These findings therefore support our hypothesis that the optimal parameter-to-compute ratio established for bpe also extends to {\textsc BLT}, or is at minimum a close approximation.

Achieving alignment with bpe scaling patterns necessitates both enhancements to model architecture and the incorporation of dynamic patching. As shown in ~\autoref{fig:scaling-compute-opt} (left), we compare models based on space-patching approaches with Llama 3. To estimate SpaceByte \citep{slagle2024spacebyte}, we employ {\textsc BLT} space-patching configured without n-gram embeddings or cross-attention mechanisms. While SpaceByte demonstrates gains over Megabyte, a substantial performance gap relative to Llama 3 persists. In \autoref{fig:scaling-compute-opt} (right), we depict the combined contributions of novel architectural design and dynamic patching. At scale, {\textsc BLT} models achieve performance competitive with the latest tokenizer-based systems, such as Llama 3.

Furthermore, we find that the selected tokenizer significantly impacts the outcomes of tokenizer-based models. In particular, those trained with the Llama-3 tokenizer surpass counterparts trained with the Llama-2 tokenizer, even when utilizing identical datasets.

In conclusion, the {\textsc BLT} architecture occupies a middle ground between Llama 2 and 3 when evaluated with much larger patch sizes. Notably, the bpe tokenizers employed by Llama 2 and 3 yield average token lengths of 3.7 and 4.4 bytes, respectively. By contrast, {\textsc BLT} is able to demonstrate analogous scaling behavior even when the average patch size is increased to 6 or 8 bytes. Since the number of inference flops is inversely related to the average patch size, employing 8-byte patches can result in almost a 50\% reduction in inference flops. Moreover, increasing the patch size tends to enhance model performance as both the scale of the model and the dataset grow. For instance, although {\textsc BLT} with an 8-byte patch size initially performs considerably worse than the bpe-based Llama 2 at the 1B parameter scale, it ultimately surpasses the bpe approach at the 7B scale. This pattern indicates that larger patch sizes may provide further benefits at even greater model and dataset sizes, and that it may be viable to explore still larger patches as computational resources increase.

\subsection{Beyond Compute Optimal Task Evaluations}
\label{section:evals}
To more thoroughly investigate scaling behaviors, we extend the training of an 8B {\textsc BLT} model past the compute-optimal threshold, using the {\textsc BLT}-1T dataset, which is both larger and of higher quality. Model performance is subsequently assessed using a range of established benchmarks covering classification and generation tasks. The selected benchmarks span common sense reasoning, world knowledge, and code generation capabilities:
\paragraph{Classification tasks} comprise ARC-Easy (0-shot)~\citep{Eval_ARC}, Arc-Challenge (0-shot)~\citep{Eval_ARC}, HellaSwag (0-shot)~\citep{Eval_hellaswag}, PIQA (0-shot)~\citep{Eval_piqa}, and MMLU (5-shot)~\citep{hendrycks2020measuring}. We utilize prompt-scoring by computing the likelihood assigned to different answer choices and present the mean accuracy achieved.
\paragraph{Coding related generation tasks:} To evaluate the model's proficiency in generating Python code, we provide pass@1 results on MBPP (3-shot)~\citep{Eval_MBPP} and HumanEval (0-shot)~\citep{Eval_HumanEval}.

In \autoref{tab:evals}, we present a comparative analysis of three models trained using the {\textsc BLT}-1T dataset: a model based on the Llama 3 tokenizer with byte pair encoding (BPE),\footnote{We opt for the Llama 3 tokenizer, featuring a 128k vocabulary, as it outperforms Llama 2’s 32k vocabulary.} alongside two distinct configurations of the {\textsc BLT} architecture. One configuration applies the space-patching method ({\textsc BLT}-Space), while the other employs an entropy-driven patching approach ({\textsc BLT}-Entropy). The entropy-based model is further constrained to approximate monotonicity and resets its context at new line boundaries, as elaborated in~\autoref{section:entropy-context}. All models are allocated a similar compute (FLOP) budget for training. Notably, for the {\textsc BLT}-Entropy variant, we adjust the entropy threshold during inference from 0.6 to 0.1, a modification that enhances performance on tasks albeit increasing the number of inference steps required.

The {\textsc BLT}-Entropy model achieves superior results compared to the Llama 3 model on 4 out of 7 evaluated tasks, despite both models being trained with an equivalent byte count. This performance gain can likely be attributed to (1) the more efficient utilization of training compute enabled by dynamic patching and (2) the explicit modeling at the byte level rather than at the token level.

Conversely, {\textsc BLT}-Space demonstrates inferior performance relative to the Llama 3 tokenizer on all tasks except one. Nevertheless, it provides notable efficiency gains by lowering inference flops, primarily due to its higher mean patch size of 6 bytes. In contrast, both bpe and entropy-patching approaches yield average patch sizes close to 4.5 bytes on the same mixture of training data. Given a fixed training compute budget, the model with larger patches is able to process 30\% more data compared to the other two, which may lead {\textsc BLT} to diverge further from the compute-optimal regime.

\subsection{Patches Scale Better Than Tokens} {\textsc BLT} models enable concurrent scaling of both model and patch dimensions without exceeding a predetermined computational budget for training and inference or altering the quantity of training samples. Patch-based architectures possess the distinctive capacity to expand patch size arbitrarily, thereby circumventing the efficiency constraints that limit traditional fixed-vocabulary token-based approaches, as detailed in Section~\ref{section:bpe-patch}. By increasing the length of patches, computational demands decrease; the resulting resources can then be redirected to enlarge the global latent transformer, since it processes input less frequently.

We perform a scaling study at fixed inference cost to investigate whether increasing model size while reducing the number of processing steps over larger input patches yields superior performance compared to employing smaller models with more steps. Using models with 400m and 3.6B parameters tokenized with Llama 2, we identify computationally comparable models employing the Llama 3 tokenizer, as well as {\textsc BLT}-Entropy variants configured with mean patch sizes of 6 and 8 bytes in the training data mixture (refer to~\autoref{tab:fixed-inf-params} for specific model configurations). For the models utilizing an 8-byte patch size, we adopt 3 encoder layers instead of a single one. Each of these models is trained under several different training compute (flop) constraints.

As illustrated in \autoref{fig:fixed_inference_scaling}, {\textsc BLT} architectures demonstrate more favorable scaling properties compared to tokenization-based models across both inference flop categories. Initially, BPE models show superior performance when the training budget is limited, but their advantage is short-lived; as the training budget increases beyond the compute-optimal point, {\textsc BLT} models begin to outperform them. From a practical standpoint, allocating greater resources to one-time pretraining can yield models that exhibit enhanced performance for a predetermined inference compute allocation. An illustrative case of this approach is found in the 8B model category, exemplified by Llama 3.1, which was pretrained on a data volume roughly two orders of magnitude higher than what would be considered optimal for its model size under compute-constrained settings.

The threshold at which {\textsc BLT} surpasses token-based models in performance has moved incrementally closer to the compute-optimal configuration as model size increases, shifting from three times to approximately 2.5 times the optimal compute allocation in larger flop-class models. Likewise, the model employing an 8-patch size demonstrates a sharper scaling trajectory within the higher flop class, overtaking alternative models at an earlier point. As elaborated in Section~\ref{section:parameter-matched-scaling}, models with larger patch sizes tend to narrow the performance gap with BPE models when evaluated at greater model capacities. This phenomenon can be partly attributed to the reduced proportion of overall computation—measured in flops—devoted to the byte-level Encoder and Decoder, which appear to scale at a slower rate than the Latent Transformer component. Specifically, when total parameters are increased twenty-fold from 400M to 8B, the number of local parameters within {\textsc BLT} only sees an approximate twofold rise. This distinction is critical, since an increase in patch size primarily impacts the flops used by the patch Latent Transformer, leaving the computational demands of byte-level modules largely unchanged. Consequently, the ratio of {\textsc BLT}-Entropy ps=8 to Llama 2 model size increased marginally from 1.6x to 1.7x as the model scale was expanded.

To conclude, our investigation into patch-length scaling reveals that the {\textsc BLT} patch-based framework exhibits superior scaling behavior when patch size and model capacity are increased together. Notably, this positive scaling appears to hold and potentially becomes more pronounced as models grow larger.

\section{Byte Modeling Improves Robustness} Additionally, we assess the robustness of {\textsc BLT} in comparison to token-based architectures that do not leverage explicit byte-level inputs, and introduce a method to incorporate byte-level information into pretrained token-based models.
\label{sec:robustness}

\subsection{Character-Level Tasks}

One of the initial incentives for developing byte-level models was their inherent resilience to noise at the byte level in input data, as well as their ability to recognize and utilize the subcomponents within tokens—a challenge for existing models that rely on tokenizers. To systematically assess these capabilities, we conduct supplementary tests using benchmarks that probe for both noise robustness and character-level understanding across various languages, incorporating English and multilingual contexts, as well as including tasks focused on digits and phonemes. The outcomes of these assessments are summarized in Table~\ref{tab:char_tasks}.

\paragraph{Noisy Data} To assess the robustness of {\textsc BLT} compared to tokenizer-based models, we generate noised variants of the benchmark classification datasets outlined in Section~\ref{section:evals}. Five character-level perturbation methods are used to diversify the text: (a) \textit{AntSpeak}: Converts all text into uppercase letters separated by spaces; (b) \textit{Drop}: Randomly deletes 10\% of characters from the sequence; (c) \textit{RandomCase}: Randomly toggles the casing, assigning uppercase to 50\% of the characters and lowercase to the remaining 50\%; (d) \textit{Repeat}: Randomly selects 20\% of characters to repeat up to four times each; and (e) \textit{UpperCase}: Uniformly changes every character in the string to uppercase. Each noise type is independently applied to the prompt, the completion, or both, and their individual effects are evaluated. The mean performance across these scenarios is then reported. As illustrated in Table \ref{tab:char_tasks}, evaluation on noised HellaSwag~\citep{Eval_hellaswag} indicates that {\textsc BLT} consistently demonstrates greater resilience than tokenizer-based models, achieving an average improvement of 8 points compared to models trained on the identical data. {\textsc BLT} also surpasses the Llama 3.1 model, despite the latter being trained on a substantially larger corpus.

\paragraph{Phonology - Grapheme-to-Phoneme (G2P)} We evaluate {\textsc BLT}'s performance in converting sequences of graphemes—letters representing words—into their corresponding phonemic (pronunciation) transcriptions. Table \ref{tab:char_tasks} displays G2P outcomes under a 5-shot experimental setup, utilizing Phonology Bench~\citep{suvarna2024phonologybench}. The results demonstrate that {\textsc BLT} surpasses the baseline model using the Llama 3 1T tokenizer for this task.

\paragraph{CUTE}
To evaluate character-level comprehension, we apply {\textsc BLT} to the CUTE benchmark~\citep{edman2024cute}. This benchmark features multiple tasks grouped into three main categories: comprehension of composition, orthographic similarity, and manipulation of character sequences. Notably, CUTE is particularly challenging for most models that rely on tokenizers, as these systems tend to know the spelling of their tokens but exhibit difficulty in using this knowledge for text manipulation. As reported in Table~\ref{tab:char_tasks}, {\textsc BLT}-Entropy surpasses both BPE Llama 3 models by over 25 points on the benchmark. The model demonstrates remarkable accuracy on character manipulation tasks, achieving 99.9\% on both spelling-related tasks. These substantial gains are even more striking considering {\textsc BLT} is trained on only one-sixteenth the data of Llama 3.1, underscoring the inherent difficulty BPE models face in acquiring robust character-level representations. Figure \ref{fig:cute_examples} presents examples where the Llama 3 tokenizer-based model fails, yet {\textsc BLT} succeeds. The only tasks where BPE models hold an advantage are word deletion and insertion. While these forms of word-level manipulation are less intuitive for byte-level approaches, the performance difference is not substantial, and constructing word-level operations from characters may be more feasible than inferring character-level behaviors from token-based representations. For all evaluations, we maintain consistent experimental protocols and utilize the original Huggingface prompts. It is possible that BPE-based models could see improvement with more tailored prompt engineering.

\paragraph{Low Resource Machine Translation} We assess the performance of {\textsc BLT} on both translation into and from English, focusing on six major language families and a set of twenty-one low-resource languages, each utilizing different scripts, as provided in the FLORES-101 benchmark~\citep{goyal2022flores}. SentencePiece BLEU scores, presented in Table~\ref{tab:flores}, summarize these evaluations. Our findings indicate that {\textsc BLT} yields consistently higher results than a baseline model utilizing the Llama 3 tokenizer, with an overall improvement of 2 BLEU points for English-targeted translations and a 0.5-point gain for translations originating from English. For widely-used language pairs, {\textsc BLT} matches or slightly surpasses Llama 3’s performance. Notably, for many language pairs belonging to low-resource families, {\textsc BLT} clearly exceeds the capabilities of Llama 3, highlighting the utility of byte-level modeling for robust handling of diverse and uncommon byte sequences.

\subsection{Training {\textsc BLT} via Llama 3 Initialization}
\label{sec:distillation}
We investigate a strategy wherein {\textsc BLT} models benefit from established pre-trained tokenizer-based models to enhance both training efficiency and convergence speed. Specifically, we initialize the global transformer parameters of {\textsc BLT} with those from a pre-trained Llama 3.1 model. Thereafter, the weights of the global transformer are fine-tuned at a learning rate set to one-tenth that of the local encoder and decoder, matching the optimal step count for Llama 3. We report comparative results against a standard {\textsc BLT} baseline in Table~\ref{tab:distill}. The findings indicate that the {\textsc BLT} model initialized from Llama 3.1 surpasses both the Llama 3 baseline and the standard {\textsc BLT}—all trained under equivalent FLOPs constraints. Furthermore, despite being trained on a much smaller dataset than the {\textsc BLT}-Entropy variant (trained on 1T tokens, cf. Table~\ref{tab:evals}), the {\textsc BLT} from Llama 3.1 demonstrates superior results on the MMLU benchmark, implying this methodology is effective in markedly decreasing the training FLOPs required.

This approach can be interpreted as a method for converting models that rely on tokenizers into ones that operate without them, thereby transforming a pre-trained LLaMA 3.1 model into a {\textsc BLT} architecture. For a thorough assessment, we present the results for the original LLaMA 3.1—trained on 15T tokens—in Table \ref{tab:distill} and directly compare its performance with that of the {\textsc BLT} derived from LLaMA 3. While our adapted model shows only a minor decrease in performance on MMLU and HumanEval, its performance is noticeably reduced on several other benchmarks. These findings indicate that additional efforts are necessary to make full use of the pre-trained model’s capabilities, especially by refining data mixing strategies and tuning other critical hyperparameters.

\section{Ablations and Discussion}
\label{section:ablations}
Here, we present a series of ablation studies that rationalize the design decisions behind the architecture of {\textsc BLT}, as well as the patching strategy and choice of hyper-parameters for the {\textsc BLT} 8B parameter model used with the {\textsc BLT}-1T dataset.

\paragraph{Entropy Model Hyper-parameters} To evaluate how scaling the size of the entropy model and altering the context window length affect overall model performance, we conduct training on byte-level entropy transformer models ranging in size from 1 million to 100 million parameters, utilizing context window lengths that vary between 64 and 512. We report the results by plotting the relationship between bits per byte (bpb) and training flops for the $400m$ and $1b$ {\textsc BLT} models, each trained with the Llama-2 dataset, as depicted in \autoref{fig:entropy_ablation}. The findings indicate that increasing both model size and context length leads to improved scaling behavior, although returns plateau once the model size surpasses 50 million parameters.

\paragraph{Types of Patching}

We evaluate the effectiveness of the four patching strategies described in Section~\ref{section:patching}: (1) Strided Patching employing strides of 4 and 6, (2) Whitespace-based Patching, (3) BPE Tokenizer patching utilizing the Llama 3 tokenizer, and (4) Entropy-based patching leveraging a lightweight byte-level language model.

Although dynamic patching decreases the practical length of sequences, we ensure that sequence length is standardized across all patching approaches to preserve comparable context windows. Across both training and inference, every model is exposed to an equivalent expected byte count per sequence, thereby eliminating potential confounds stemming from access to longer contexts. The findings of these ablation studies are illustrated in Figure~\ref{fig:scaling-compute-opt}. Every alternative patching strategy surpasses the static patching baseline, with space patching showing performance nearly on par with dynamic entropy based patching.

\autoref{tab:evals_ablation} displays comparative benchmark results for {\textsc BLT} models utilizing tokenizer-based approaches, space patching, and entropy-based patching, all trained on the Llama 2 dataset for the empirically determined optimal number of training steps~\citep{dubey2024llama}. While space patching offers a more straightforward alternative by sidestepping the need for dynamic entropy model computation during training, our findings reveal that the performance improvements seen with entropy-based patching in the context of scaling laws~(Section~\ref{section:scaling}) also extend to practical downstream benchmarks.\footnote{Space patching results were obtained from prior experiments without cross-attention; nonetheless, the overall patterns persist when cross-attention is incorporated.}

\paragraph{Cross-Attention} 
In~\autoref{tab:crossattn_ablations_1b}, we investigate the effects of adding cross-attention at different stages within both the encoder and decoder components of {\textsc BLT}. Specifically, for encoder cross-attention, we compare several methods for initializing the queries: 1) utilizing a shared, learned embedding across all global states, 2) applying a hash embedding to the bytes within the patch, and 3) using a pooled version of the encoder’s hidden state representation corresponding to the patch bytes at the specified encoder layer.

Our results indicate that incorporating cross-attention within the \textit{decoder} yields the greatest performance benefits. When implemented in the encoder, cross-attention offers marginal gains, but these are observed exclusively when queries are initialized via pooling. Furthermore, cross-attention proves especially advantageous on the Common-Crawl dataset and demonstrates marked improvements when applied with increased patch sizes.

\paragraph{n-gram Hash Embeddings} 

We experiment with n-gram hash embedding vocabularies set to 0, 100K, 200K, and 400K, with the corresponding outcomes reported in Table~\ref{tab:ngram_ablations_1b}. Our analysis demonstrates that hash embeddings confer benefits across all tested domains, with especially marked improvements on Wikipedia and Github (yielding a 0.04 bpb difference compared to a 0.01 bpb difference after 15k training steps at the 8B parameter scale). At the 8B model size, reducing the number of hashes from 500K to 300K impacts results by roughly 0.001 bpb after 15k steps. These observations suggest that incorporating hashes is crucial for aligning the performance of {\textsc BLT} models with that of those employing tokenizers, yet enhancements plateau beyond 300K hashes. Furthermore, evidence suggests that the improvements from hashing and those from cross-attention mechanisms are mostly additive, with each method benefiting distinct datasets.

\paragraph{Local Model Hyperparamaters} Table \ref{tab:local_layers} presents an ablation of different choices regarding the layer count in both the local encoder and decoder. For hash n-gram embeddings, {\textsc BLT} achieves strong results using a minimal encoder—specifically, a single-layer configuration—while performance improves with a more complex, multi-layer decoder.

\section{Related Work}
\label{section:related_work}

\textbf{Character-Level RNNs:} Since the advent of neural approaches, Character Language Modeling has garnered significant interest due to its natural capacity to handle out-of-vocabulary words, thereby eliminating the necessity for explicit back-off strategies~\citep{sutskever2011generating,mikolov2012subword,graves2013generating}. \cite{kim2016character} introduced a system that exclusively processes characters on the input side via convolutional and highway layers preceding LSTM-based RNNs; their model achieved competitive results with the then state-of-the-art RNN language models on English, while also surpassing them on languages with complex morphology—highlighting an important benefit of character-level LLMs. Similarly, \cite{kenter2018byte} demonstrated superior machine comprehension performance using byte-level LSTM architectures over word-level models, particularly for morphologically rich languages such as Turkish and Russian. Extending this idea, \cite{zhang2015character} applied convolutional character-based models to classification problems and found they could outperform word-based alternatives in certain contexts. In a further advancement, \citet{chung2019hierarchical} developed hierarchical LSTM frameworks with boundary detectors at each tier, enabling the models to uncover latent structural hierarchies in text, thus enhancing results in character-level language modeling. Rather than relying on attention mechanisms, ByteNet by \cite{kalchbrenner2016neural} leverages convolutional layers over character sequences for the task of machine translation.

\textbf{Character-Level Transformers:} Transformer architectures that utilize attention mechanisms~\citep{vaswani2017attention} in conjunction with subword tokenization methods~\citep{sennrich-etal-2016-neural} have led to substantial advancements in neural language modeling and various benchmark evaluations. Nevertheless, adopting word or subword representations inherently imposes a particular inductive bias regarding the abstraction level at which the models function. In an effort to merge the robust capabilities of transformers with the positive early outcomes in character-level modeling, \citet{al2019character} implemented exceedingly deep transformer networks and incorporated auxiliary loss functions. This approach enabled transformer-based character language models to surpass previous LSTM-driven character LLMs. Despite these improvements, a notable performance disparity persisted between these models and those trained at the word level. Additionally, GPT-2~\citep{radfordgpt2} reported that byte-level language models fell short of their word-level counterparts when evaluated on extensive datasets such as the 1 billion word benchmark.

\cite{Choe2019BridgingTG} provided evidence that transformer-based byte-level LLMs can surpass subword-level models with similar parameter counts, but achieving such performance requires substantially greater computational resources and longer training times. In a similar vein, \cite{el2020characterbert} introduced a BERT variant known as CharFormer, which constructs word representations by convolving over character embeddings; although this model demonstrated effectiveness in the medical domain, its training was significantly more computationally demanding. The CANINE model proposed by \cite{clark2022canine} employs 150 million parameters and processes input at the character level. At its core, CANINE incorporates a deep transformer architecture, akin to the global model described here, and leverages both a local transformer and strided convolution layers to efficiently reduce the input sequence length. Notably, CANINE achieves superior results over token-based encoder models, such as mBERT, on a range of multilingual downstream tasks. The ByT5 model~\citep{xue2022byt5} explored the efficacy of encoder-decoder architectures at the byte level that entirely omit patching strategies. While ByT5 exhibited enhanced resilience to noise and performed competitively with models reliant on tokenization—especially when trained with four times less data—the necessity of computing self-attention over every individual byte imposed considerable computational overhead. One intrinsic challenge with byte-level modeling is the elongated input sequence, which exacerbates the difficulty of scaling such models efficiently. The emergence of the Mamba Architecture~\citep{gu2023mamba}, capable of sustaining a fixed-size memory state even across long input contexts, enabled \cite{wang2024mambabyte} to train a byte-level Mamba-based model without employing patching. Their approach achieved better performance than byte-level transformers—evaluated by bits-per-byte—in flop-constrained scenarios for models at the 350M parameter scale across several datasets.

\textbf{Patching-based approaches:} Utilizing patching strategies can significantly reduce the computational cost (in terms of flops) for byte-level LLMs without necessarily compromising model effectiveness. Various studies have reported promising outcomes at modest model sizes and with limited training data. \cite{nawrot-etal-2022-hierarchical} introduce static patching through downsampling and upsampling, resulting in the hourglass transformer, which achieves superior performance relative to other byte-level models at the 150M parameter scale. Subsequently, \cite{nawrot-etal-2023-efficient} extend these results with dynamic patching techniques, such as a trainable end-to-end boundary predictor, a boundary-predictor supervised by specific tokenizers, and an entropy-guided patching method analogous to {\textsc BLT}. Their results demonstrate that these dynamic strategies can outperform standard transformers on language modeling benchmarks, particularly at the 40M parameter scale and using around 400M training tokens. Further, \cite{lester2024training} explore sequence compression via arithmetic coding, surpassing the compression levels attainable with BPE, and introduce an equal-info windows method. While this approach yields better results than byte-level baselines on language modeling tasks, it does not yet exceed the performance of subword-based models.

Our research is heavily influenced by, and most similar to, MegaByte~\citep{yu2023megabyte}, a decoder-only causal LLM architecture that employs fixed, static patching combined with concatenation of representations for converting bytes to patches. This approach is followed by a local model in the decoder that maps patches back to bytes. The authors show that MegaByte is able to achieve performance parity with tokenizer-based models at the 1B parameter scale when trained on a dataset comprising approximately 400B bytes. In all our experiments, we include MegaByte as a baseline, consistently observing that static patching falls short compared to the most recent state-of-the-art tokenizer-based models that are optimally trained with equivalent computational resources. We further illustrate how {\textsc BLT} narrows this performance gap. Similarly, \cite{slagle2024spacebyte} observe these limitations with MegaByte, proposing enhancements by adapting the patching strategy to use whitespaces and other space-like bytes as patch boundaries, and by incorporating a local encoder. Their findings suggest performance gains over token-based transformer models in certain domains (such as Github and arXiv) at the 1B parameter scale in compute-constrained regimes. We have also replicated experiments with this model, and our results indicate that additional architectural advances are required to enable byte-level models to scale further and attain parity with state-of-the-art token-based systems like Llama 3.

\section{Limitations and Future Work}

For the selection of model architectures in this study, we adopt training regimes matching the optimal step counts as specified for Llama~3~\citep{dubey2024llama}. It should be noted, however, that the original scaling laws were derived using BPE-level transformers, which means that applying them directly to {\textsc BLT} might not yield optimal ratios between data and parameter count. Investigating scaling laws specifically tailored to {\textsc BLT} remains an avenue for future research and could potentially reveal even more advantageous scaling properties for this model. Furthermore, our current experiments predominantly cover scales up to 1B parameters; architectural optima may shift when models are scaled to 8B parameters or more, suggesting that there is scope for further performance enhancements at these larger scales.

Current transformer libraries and codebases are optimized for performance with tokenizer-based transformer models. Although we provide theoretically matched FLOP comparisons and utilize efficient solutions like FlexAttention for non-standard transformer layers, our implementations may still lag behind tokenizer-based models in wall-clock speed and could potentially see improvements through additional optimization efforts.

Although {\textsc BLT} relies on an independently trained entropy model for its patching mechanism, exploring end-to-end training of the patching model itself presents a promising avenue for future investigation. Section \ref{sec:distillation} outlines preliminary experiments that demonstrate potential for ``byte-ifying'' large-scale tokenizer-based models like Llama 3—models originally trained on datasets exceeding 10T tokens—by adopting a strategy of weight initialization and freezing for the global transformer. Continued research in this area could yield approaches that not only preserve the strengths of bytefying but also potentially surpass the performance of such tokenizer-based models, eliminating the necessity for full retraining from scratch.

\section{Conclusion}
\label{section:conclusion}

In this work, we introduce the Byte Latent Transformer (\textbf{BLT}), an architectural innovation that challenges the traditional reliance on fixed-vocabulary tokenization within large language models. {\textsc BLT} implements an adaptive, learnable mechanism for aggregating bytes into patches, allowing computational effort to be distributed according to the complexity of the input data. This leads to marked gains in computational efficiency and model robustness. Through a thorough scaling analysis, we show that {\textsc BLT} attains comparable performance to token-based systems such as Llama 3 at parameter sizes up to 8B and data scales up to 4T bytes, with the capacity to reduce inference flops by as much as 50\% in exchange for modest decreases in standard evaluation scores. Additionally, {\textsc BLT} facilitates novel scaling strategies, making it feasible to expand both the model size and patch dimension concurrently within a given inference cost. This property proves advantageous for compute-constrained environments that are typical in real-world applications. By processing data directly at the byte level, {\textsc BLT} further enhances the model’s capabilities in dealing with less frequent data patterns, conferring greater resilience to input noise and promoting richer representations of sub-word information. Collectively, these findings suggest that {\textsc BLT} offers a compelling, robust, and adaptable alternative to conventional tokenization-based language modeling frameworks.

\end{document}